\documentclass[10pt]{article}

\usepackage[utf8]{inputenc}

\usepackage[T1]{fontenc}

\usepackage{amsmath}

\usepackage{amsfonts}

\usepackage{amssymb}

\usepackage[version=4]{mhchem}

\usepackage{stmaryrd}

\usepackage{bbold}

\usepackage{graphicx}

\usepackage[export]{adjustbox}

\graphicspath{ {./images/} }



\begin{document}

здесь $V(t)$ - нормированная в нуле фундаментальная матрица системы (1). Отсюда следует экспоненциальная дихотомия системы (1), если $\operatorname{Re} \alpha_{j} \neq 0$ ни для какого $j=1, \ldots, l$.



JIum. [1] ' F l o q u e t G., "Ann. sci. École norm. supér.", 1883 , t. 12, № 2 , p. 47-88; [2] Л я ІІ у н о в А. М., Собр. соч., т. 2, М. - Л., 1956, с. 7-263; [3] Д е м и д о в и ч'Б. П., Лекции по математической теории устойчивости, М., 1967; [4] Я к уб о в и ч В. А., С т а р ж и н с к и й В. М., Јинейные дифференциальные уравнения с периодическими коэффициентами и их приложения, М., 1972; [5] М а с с е p а X. - Л., III е фіпр е р X. - X., Јинейные дифференциальные уравнения и функфе p X. - Х., Линейные дифференциальные уравнения и фу финг и н Н. П., Линейные системы обыкновенных дифференциаль\begin{tabular}{l}

ных уравнений с периодическими и квазипериодическими козф- \\

Ю. В. Комленко \\

\hline

\end{tabular}



ФЛОКЕ - ЛЯПУНОВА ТЕОРЕМА - см. Флоке теория.



ФOКА ПРОСТРАНСТВО, ф о к о в с к о е п р о с тр а н с т в о,- в простейшем и чаще всего употребляемом случае - гильбертово пространство, состоящее из бесконечных последовательностей вида



\[

F=\left\{f_{0}, f_{1}, \ldots, f_{n}, \ldots\right\}

\]



где



\[

f_{0} \in \mathbb{C}, f_{1} \in L_{2}(\mathbb{R} v, d v x), f_{n} \in L_{2}^{S}\left((\mathbb{R} v)^{n},\left(d^{v} x\right)^{n}\right),

\]



или



$f_{n} \in L_{2}^{a}\left(\left(\mathbb{R}^{v}\right)^{n},\left(d^{v} x\right)^{n}\right), n=2,3, \ldots, v=1,2, \ldots$, причем



\[

L_{2}^{S}\left(\left(\mathbb{R}^{v}\right)^{n},\left(d^{v} x\right)^{n}\right) \text { или } L_{2}^{a}\left(\left(\mathbb{R}^{v}\right)^{n},\left(d^{v} x\right)^{n}\right)

\]



означает гильбертово пространство симметрич. (соответственно антисимметрич.) функций от $n$ переменных $x_{1}, \ldots, x_{n} \in \mathbb{R} v, n=2,3, \ldots$ Скалярное произведение двух последовательностей $F$ и $G$ вида (1) равно



\[

(F, G)=f_{0} \bar{g}_{0}+\sum_{n=1}^{\infty}\left(f_{n}, g_{n}\right)_{L_{2}}\left(\left(\mathbb{R}^{v}\right)^{n},\left(d^{v} x\right)^{n}\right) .

\]



В случае когда последовательности $F$ состоят из симметрич. функций, говорят о с и м м е т р и ч е с к о м (или б о 3 о н н о м) Ф. П., а в случае последовательностей антисимметрич. функций - Ф. п. наз. а н т и-



\includegraphics[max width=\textwidth, center]{2023_07_09_db5decd134ecae7502beg-1}

В таком простейшем случае Ф. п. были впервые введены В. А. Фоком [1].



В общем случае произвольного гильбертова пространства $H$ Ф.п. $\Gamma^{S^{(}}(H)$ (или $\Gamma^{a}(H)$ ), построенным над $H$, наз. симметризованную (или антисимметризованную) тензорную экспоненту пространства $H$, т. е. пространства



\[

\Gamma^{\alpha}(H) \equiv \operatorname{Exp}_{\alpha} H=\bigoplus \sum_{n=0}^{\infty}(H \otimes)_{\alpha}, \quad \alpha=S, a,

\]



где знак $\bigoplus$ означает прямую ортогональную сумму гильбертовых пространств, $\left(H^{\infty 0}\right)_{\alpha}=\mathbb{C}^{1},\left(H^{\left(\otimes^{1}\right.}\right)_{\alpha}=H$, a $\left.\left(H^{\bigotimes}\right)^{n}\right)_{\alpha}, n>1,-$ симметризованную при $\alpha=S$ или антисимметризованную $(\alpha=a) n$-ю тензорную степень пространства $H$. В случае $H=L_{2}\left(\mathbb{R} v, d^{v} x\right)$ оиределение (2) эквивалентно определению Ф. П., приведенному в начале статьи, если отождествить гространства $L_{2}^{\alpha}\left((\mathbb{R} v)^{n},\left(d^{v} x\right)\right)^{n} \mathbf{u}\left(L_{2}\left(\mathbb{R} v, d^{v} x\right)\right) \underset{\alpha}{\bigotimes}{ }^{n}$ так, что тензорному произведению



\[

\left(f_{1} \otimes \cdots \otimes f_{n}\right)_{\alpha} \in\left(L_{2}\left(\mathbb{R}^{v}, d^{v} x\right)\right)_{\alpha}^{\otimes}{ }^{n}

\]



последовательности функций



\[

f_{1}, \ldots, f_{n} \in L_{2}\left(\mathbb{R}^{v}, d^{v} x\right)

\]



соответствует функция



\[

\frac{1}{\sqrt{n !}} \sum_{\sigma}( \pm 1)^{\operatorname{sign} \sigma} \prod_{i=1}^{n} f\left(x_{\sigma(i)}\right) \in L_{2}^{\alpha}\left(\left(\mathbb{R}^{v}\right)^{n},\left(d^{v} x\right)^{n}\right)

\]



где суммирование происходит по всем перестановкам $\sigma$ индексов $1,2, \ldots, n$, sign $\sigma$ - четность перестановки $\sigma$, а знак +1 или -1 в выражении (3) соответствует симметрич. или антисимметрич. слутаю.



В квантовой механике Ф. п. $\Gamma^{S}(H)$ или $\Gamma^{a}(H)$ служат пространствами состояний квантовомеханич. системы, состоящей из произвольного (но конечного) числа одинаковых частиц таких, что пространством состояний каждой отдельной частицы является пространство $H$. При этом в зависимости от того, каким из Ф. П.- симметрическим $\Gamma^{S}(H)$ или антисимметрическим $\Gamma^{u}(H)$ описывается әта система-сами частицы наз. б о з о н а м іл иліг соответственно ф е р м и о н а ми. Для любого $n=1,2, \ldots$ подпространство $\Gamma_{n}^{\alpha}(H) \equiv$ $\equiv\left(H^{\left(\otimes^{n}\right.}\right)_{\alpha} \subset \Gamma^{\alpha}(H), \quad \alpha=S, a, \quad$ наз. $n$-ч а с т и ч ны м п о д п р о т р анст вом: его векторы опісывают те состояния, в к-рых имеется ровно $n$ частіц; единичный вектор $\Omega \in\left(H^{\bigotimes}{ }^{0}\right)_{\alpha} \subset \Gamma^{\alpha}(H), \alpha=S, a$ (в записи (1): $\Omega=\{1,0,0, \ldots, 0, \ldots\})$, наз. в а г у умны м в ектором, описывает состояние системы, в к-ром нет нІІ одной тастицы.



При изучении линейных операторов, действующих в Ф. п. $\Gamma^{S}(H)$ и $\Gamma^{a}(H)$, часто применяется специальный формализм, наз. ме то дом в то р и ч но го к в а н т о в а в я. Он основан на введении в каждом из пространств $\Gamma^{\alpha}(H), \alpha=S, a$, двух семейств линейных операторов: т. н. о п е р а т о р о в уни ч то жен ил я $\left\{a_{\alpha}(f), f \in H\right\}, \alpha=S, a$, и семейства сопряженных к ним операторов $\left\{a_{\alpha}^{*}(f) f \in H\right\}$, наз. о п е р а т о р а м и р ож д е н и. Операторы уничтожения задаются как замыкания операторов, действующих на векторы



\[

\left(f_{1} \otimes \cdots \otimes f_{n}\right) \propto \in \Gamma^{\alpha}(H), \alpha=S, a,

\]



где $\left(f_{1} \otimes . . \otimes f_{n}\right)_{\alpha}$ симметризованные (при $\alpha=S$ ) или антисимметризованные $(\alpha=a)$ тензорные произведения последовательностей векторов $f_{1}, \ldots, f_{n} \in H, n=1,2, \ldots$, по формулам



\[

\begin{gathered}

a_{\alpha}(f)\left(f_{1} \otimes \ldots \otimes f_{n}\right)_{\alpha}=\sum_{i=1}^{n}(-1)^{g_{\alpha}(i)}\left(f_{i}, f\right) \times \\

\times\left(f_{1} \otimes \ldots \otimes f_{i-1} \otimes f_{i+1} \otimes \cdots \otimes\right. \\

\left.\cdots f_{n}\right)_{\alpha}, \\

\alpha=S, a, a_{\alpha}(f) \Omega=0,

\end{gathered}

\]



где $g_{S}(i)=0$ и $g_{a}(i)=i-1$. Операторы же рождения $a_{\alpha}^{*}(f)$ действуют на векторы (3) по формулам



\[

\begin{gathered}

a_{\alpha}^{*}(f)\left(f_{1} \otimes \ldots \otimes f_{n}\right)_{\alpha}=\left(f \otimes f_{1} \otimes \ldots \otimes f_{n}\right), \\

a_{\alpha}^{*}(f) \Omega=f .

\end{gathered}

\]



При этом для любого $f \in H, a_{\alpha}(f): \Gamma_{n}^{\alpha}(H) \rightarrow \Gamma_{n-1}^{\alpha}(H)$, $n=1,2, \ldots$, а $a_{\alpha}^{*}(f): \Gamma_{n}^{\alpha}(H) \longrightarrow \Gamma_{n+1}^{\alpha}(H), n=0,1,2, \ldots$, т. е. состояние физич. системы с $n$ частицами операторами уничтожения $a_{\alpha}(f)$ переводится в состояние с $(n-1)$-й частицей, а операторами рождения $a_{\alpha}^{*}(f)-$ в состояние с $(n+1)$-й частицей. Операторы рождения и униттожения оказываются во многих слутаях удобной системой «образующих» в совокупности всех операторов (ограниченных и неограниченных), действующих в Ф.п. ПЈредсґавление таких операторов в виде суммы (конечной илі бесконечной) операторов вида $a_{\alpha}^{*}\left(f_{1}\right) \ldots a_{\alpha}^{*}\left(f_{n}\right) a_{\alpha}\left(g_{1}\right) \ldots a_{\alpha}\left(g_{m}\right)$, $\left(f_{1}, \ldots, f_{n}, g_{1}, \ldots, g_{m} \in H, n, m=0,1,2, \ldots\right)$,



т. н. но рма льна я фо р ма о пе р а то ра,и основанные на таком представлении способы действия с операторами (вычисление функций от них, приведение операторов $\kappa$ какому-нибудь «простейшему" виду, различные приемы аппроксимации и т. д.) и составляют содержание упомянутого выше формализма вторичного квантования (см. [2]).





\end{document}